% The NeurIPS abstract environment supplies the unnumbered Abstract heading.
\begin{abstract}
Can an outline help a language model withhold a secret while writing a story?
We compare secret-blind outlines with a plain secrecy instruction, exactly length-matched irrelevant context, and a decoy-word control using Gemma~3~12B.
Across 90 word-secret cases and 60 synthetic plot cases, we generate 1,080 stories and evaluate secret-present versus no-secret pairs in both presentation orders with Qwen3~14B.
For word secrets, detection accuracy falls from 63.9\% with irrelevant context to 54.4\% with outlines, a $-9.4$ percentage-point difference (95\% CI $[-15.0,-3.9]$; Holm-adjusted $p=.0155$).
However, literal disclosures dominate the observed signal: the forbidden word appears in 27/90 baseline stories, 15/90 outline stories, and 31/90 context stories.
A secondary likelihood-based detector does not confirm any multiplicity-corrected contrast, and plot results remain inconclusive.
A fixed cross-context residual decoder achieves only 8.9\% accuracy over 15 secrets, failing the behavioral-intervention gate.
Post hoc counterfactual centering recovers secret identity at 82.2--91.1\% across baseline continuation positions, showing recoverable conditioning information without establishing causal influence.
Strong evaluator position bias and a constant duplicated beginning-of-sequence token limit generalization.
The supported finding is narrower than general secrecy: under this recorded protocol, outlines reduce exact-word violations, while nonliteral leakage and selective causal removal remain unresolved.
\end{abstract}
